% ch8_discussion.tex — 第八章 讨论
\section{讨论}

\subsection{H1讨论：ATT对Rating的预测效度}

H1的发现揭示了LLM提取的多属性态度指数与平台客观评分之间的强预测关系（有序Logistic $\beta=1.867$，$p<0.001$；OLS $R^2=0.560$），为"Text as Data"方法论在消费者行为研究中的应用提供了实证支持。这一结果与Mudambi \& Schuff（2010）\cite{mudambi2010}关于评论诊断性的研究形成呼应——该研究发现评论文本的信息深度影响消费者的评分感知，而本研究进一步表明，从文本中系统提取的多属性信念结构本身就与评分存在高度对应（$r=0.722$）。

这里需要回应一个潜在质疑：ATT与Rating的高相关是否意味着两者测量的是同一构念，从而构成"循环论证"？本研究认为不然，理由有三。其一，两者的测量机制根本不同：ATT通过LLM对非结构化文本进行多属性分解得到，Rating是评论者在行为完成后的单次整体判断，信息来源和加工过程均独立；其二，H4的增量效度检验（$\Delta R^2=0.018$，$p=0.004$）证明，在控制Rating之后ATT仍对行为意向有额外预测力，说明两者并不等价；其三，高相关本身正是聚合效度（convergent validity）的体现——若ATT对评分毫无预测力，才应怀疑LLM提取方法的有效性。因此，$r=0.722$既是方法论验证的核心证据，也为后续研究将LLM提取的多属性态度用于评分预测提供了依据。

\subsection{H2讨论：评论者经验的调节效应}

H2的发现揭示了一个在理论上具有启示意义的反向调节模式：低经验评论者的ATT--Rating一致性（简单斜率$=0.741$）反而强于高经验评论者（$0.573$），这与基于Fazio \& Zanna（1978）\cite{fazio1978}直接经验理论"经验增强态度--行为一致性"的预期相悖。

这一发现不应被解读为理论的失效，而是揭示了经验效应在在线评论情境下的特殊作用机制。本研究提出以下解释路径：

\textbf{比较性评价框架的形成。}高经验评论者（Beauty品类评论数多）在长期评论实践中形成了跨产品的参照体系。其评分行为不仅反映对当前产品属性的直接判断，还融入了与竞品、替代品和历史经验的比较性权衡。这种比较性评价框架使ATT（仅基于当前评论文本的属性信念）无法完整解释评分行为，导致两者的对应关系相对松弛。低经验评论者缺乏此类跨产品参照，其评分更直接地映射当前产品体验，ATT与Rating因此保持更高一致性。

\textbf{评分策略的保守化偏移。}Mitchell（2022）\cite{mitchell2022}的研究发现，消费者专业性水平越高，其评分行为越呈现出"保留高分区间"的策略性特征。高经验评论者可能形成了相对固定的评分锚点（如"5星仅授予真正卓越的产品"），即使多属性态度积极，也不轻易给出最高评分，造成ATT与Rating之间的系统性偏差，而这种偏差在低经验评论者中较少出现。

H2的发现对消费者行为研究具有双重理论意义：一方面，它挑战了"经验线性增强态度--行为一致性"的简单假设，表明经验的调节效应是情境依赖的；另一方面，它提示在利用LLM提取多属性态度进行评分预测时，需要将评论者经验纳入模型，对高经验群体的预测精度期望做出相应调整。

\subsection{H3讨论：ATT对BI的直接效应}

H3的发现（$\beta=0.167$，$p=0.004$）与Fishbein \& Ajzen（1975）\cite{fishbein1975}理性行为理论的核心命题一致：态度是行为意向的近端预测变量。在同时控制Rating和控制变量的条件下，ATT仍对BI保持独立预测力，说明多属性分解捕捉到了单一评分无法覆盖的信念结构信息。

需要在方法论框架下审慎解读这一结论。H3/H4分析存在两项已知局限：其一，ATT与BI均提取自同一评论文本，Harman单因子检验（76.8\%$>$50\%阈值）证实存在严重的同源方差（CMV）问题，两者的相关可能被人为放大；其二，BI存在组（$N=298$）的ATT均值（1.370）显著高于BI缺失组（0.944），$t=4.604$，$p<0.001$，选择偏差使子样本的态度分布偏向积极端，H3的效应量可能高于总体真实水平。因此，H3的结论仅适用于"明确在评论中表达行为意向的评论者"这一子群，不宜推广至全体评论者。

\subsection{H4讨论：ATT的增量效度}

H4的发现是本研究最重要的理论贡献之一：在已解释大部分BI方差的Rating（$\Delta R^2=0.314$）之上，ATT仍能提供显著增量预测力（$\Delta R^2=0.018$，$p=0.004$）。

$\Delta R^2=0.018$的绝对量虽小，但其理论意义不应以效应量的大小来衡量，而应以其所揭示的信息增量来理解。评分作为单一整体判断，压缩了消费者态度结构中大量的内部分化信息——例如，"外包装简陋但保湿效果优异"的评论者可能给出4星，其高保湿信念强度与低外观满意度相互抵消后才形成评分，而ATT的多属性分解能够保留这种属性间的结构性差异。正是这部分单一评分所"丢失"的信念结构，构成了ATT对BI的增量预测来源。这一解释与Fishbein多属性态度理论\cite{fishbein1975}的核心论点一致：相比整体态度，属性分解视角能更精确地捕捉消费者行为意向背后的信念机制。

从信息论的视角看，Rating作为单一值已将消费者态度的绝大部分信息压缩其中（Model 2 $R^2=0.336$），剩余可解释的BI方差本就十分有限；ATT能在这一狭小空间中额外捕获1.8\%的方差，恰恰说明多属性分解触及了单维评分无法覆盖的态度结构性差异。换言之，这不是ATT的预测力弱，而是Rating已经很强——ATT在高度压缩的剩余信息空间里仍能挤出显著增量，正是其独立信息价值的体现。

从方法论层面，H4的支持回应了"ATT仅是Rating代理变量"的质疑，证明多属性分解具有独立的信息价值，为在行为预测模型中同时纳入评分与多属性态度提供了依据。

\subsection{方法论讨论：LLM文本提取的可靠性}

LLM可靠性验证结果（ATT计算准确率100\%，幻觉率4.2\%，跨模型重测信度$r=0.892$）表明，在规范的提示词设计（防幻觉约束、few-shot校准、格式强制）下，LLM能够可靠地执行基于Fishbein理论的多属性信念提取任务。

值得关注的是$b_i$信念强度的ICC（0.249），低于0.70的可接受阈值。这一结果并不意味着整体方法失效——$e_i$情感极性的跨模型一致率（97.2\%）和ATT总分的重测信度（$r=0.892$）均达到良好水平，说明信念强度赋值的绝对数值存在模型间差异，但情感方向判断和加权汇总结果保持稳定。其可能的原因是不同LLM在0--3量表的数值校准上缺乏统一参照，这一局限提示未来研究可通过锚定样例（anchor examples）在提示词中统一量表使用标准。

三数据源设计（LLM文本提取/平台客观评分/全品类行为数据）是本研究在方法论上的核心贡献。这一设计将Podsakoff等（2003）\cite{podsakoff2003}提出的程序控制策略与数据来源独立性原则结合，为利用大规模平台数据进行消费者行为研究提供了可复用的CMV控制范式，具有超越本研究情境的方法论推广价值。

\subsubsection{消除CMV的进一步尝试与方法论启示}

为探究能否从根源上消除H3/H4分析的同源方差，本研究事后探索了两种将行为意向操作化为外部行为指标的可能性。

\textbf{方案一（重复购买行为）。}以评论者后续是否再次评论同一产品作为实际回购行为的代理变量。在目标产品（ASIN: B004EDYQX6）的1{,}331条评论中，仅发现7位评论者（涉及15条评论，占总样本1.1\%）有重复评论行为。为排除该比例仅为个案特殊性所致，进一步对同品类评论量最大的产品进行验证：以REVLON One-Step Volumizer（37{,}185条评论）为例，重复评论用户仅319人（0.87\%），且该比例在不同规模的产品中保持高度稳定。即便将分析范围扩大至评论量超过10万条的产品，按0.87\%的结构性比率推算，可获取的重复评论者约870人——虽在统计效力上勉强达标，但该代理指标存在三个根本性的构念效度缺陷：其一，绝大多数消费者回购后不再次评论，``未重复评论''被错误等同于``未回购''，假阴性率极高，系统性地稀释效应量；其二，同一用户对同一产品的多条评论中，相当一部分属于更新评论或补充评论而非独立购买后分别撰写，``重复评论''与``重复购买''之间的构念映射存在根本性缺陷；其三，计划行为理论中行为意向（BI）与实际行为（Behavior）属于不同理论层面的变量\cite{ajzen1991}，意向--行为之间存在已被广泛记录的鸿沟\cite{sheeran2022}，以实际回购行为替代行为意向作为效标，实质上改变了H3/H4的理论检验结构。综合以上分析，方案不可行。

\textbf{方案二（有用票数代理）。}将评论的\texttt{helpful\_vote}（有用票数）作为口碑推荐行为的客观代理。然而，89.0\%的评论未获得任何有用票（极度零膨胀分布），且\texttt{helpful\_vote}在理论上测量的是信息接收者对评论质量的感知，而非评论者自身的行为意向，存在构念上的根本错位，方案同样不可行。

上述探索的结果从反面凸显了三数据源设计的独特价值与严苛的前提条件。它进一步表明，在利用用户生成内容开展行为研究时，能够提前规划并实现测量来源的彻底分离（如H1/H2分析的设计），是一种罕见但关键的方法论优势——其价值正是通过"替代方案均行不通"这一事实而得以量化。这一探索也为未来研究指明了方向：要彻底解决态度--行为意向关系中的CMV问题，可能需要依靠电商平台提供匹配的后续购买行为数据，或采用前瞻性的纵向研究设计。
